\section{Conclusion}
\label{sec:conclusion}

We argued that autonomous index selection should be posed as a continual
reinforcement-learning problem rather than a one-shot optimization, because
production workloads drift and both offline retraining and naive online
fine-tuning fail---the former is expensive, the latter forgets. We introduced
CLDB, an index-aware DQN agent that combines drift-triggered elastic weight
consolidation, database-aware prioritized replay, and safety-constrained
exploration, all running off the query critical path. We gave a complete
mathematical formulation, a theoretical analysis of the EWC objective, the
value-network Fisher estimator, the prioritized-replay gradient, and the
consolidation cost, and we built a reproducible benchmark of conflicting workload
phases. Empirically, the agent that learns without replay forgets catastrophically
while the replay-based agents retain previously learned phases and beat static,
random and heuristic baselines.

The clearest path to a definitive study is to replace the cost model with a real
DBMS and a standard benchmark workload, add published index-advisor baselines,
increase the number of seeds and phases, and make the EWC strength adaptive. We
release the implementation and the benchmark so that these experiments can be
run and extended.

\paragraph{Reproducibility.}
All code, the simulator, the benchmark harness and the figure-generation script
are provided. The benchmark is configurable through environment variables and
seeded; the paper's figures and tables are produced directly from the exported
results file.
